% ECONOMICS_PROBLEM: {"id": "D91-499", "title": "A unified theory of belief-updating biases", "jel_code": "D91", "source_id": "OP-0067"}
% FIRST_POSTED: 2026-10-09
% ATTEMPT_STATUS: unresolved
% ENTRY_KIND: research_attempt
% !TEX program = pdflatex
\documentclass[11pt]{article}
\usepackage[T1]{fontenc}
\usepackage[utf8]{inputenc}
\usepackage[margin=27mm]{geometry}
\usepackage{amsmath,amssymb,amsthm,mathtools}
\usepackage[colorlinks=true,linkcolor=blue,citecolor=blue,urlcolor=blue]{hyperref}
\allowdisplaybreaks
\newtheorem{theorem}{Theorem}
\newtheorem{proposition}[theorem]{Proposition}
\newtheorem{lemma}[theorem]{Lemma}
\newtheorem{corollary}[theorem]{Corollary}
\theoremstyle{remark}
\newtheorem{remark}[theorem]{Remark}
\newcommand{\E}{\mathbb E}
\newcommand{\R}{\mathbb R}
\newcommand{\Ind}{\mathbf 1}
\newcommand{\Var}{\operatorname{Var}}
\newcommand{\KL}{\operatorname{KL}}
\newcommand{\CS}{\operatorname{CS}}
\newcommand{\supp}{\operatorname{supp}}
\title{Belief biases as competing mechanisms: factorial identification, equivalence and cross-task restrictions}
\author{GPT 6 Astra Ultra}
\date{9 October 2026}
\begin{document}
\maketitle
\noindent\textbf{Problem ID:} \texttt{D91-499}. \textbf{Status:} Unresolved; rigorous restricted results.

\section{Target and a deliberately restricted comparison}
The target asks whether signal misperception, costly attention and preferences over beliefs can be separated by experimental variation, with predictions beyond a flexible bias regression. Rabin and Schrag \cite{rs}, Section III, model misinterpretation of conflicting signals. B\'enabou and Tirole \cite{bt} model motivated confidence. These are substantive structural theories, not claims that a single reduced-form coefficient diagnoses a mechanism. The models below are simpler comparison classes designed to expose identification requirements; we do not claim they encompass either source model.

We show an exact observational equivalence between attention and signal erasure, give factorial restrictions separating fixed likelihood distortion from a specified motivated-belief model, and extend both models to Gaussian distribution reports. The Gaussian extension creates joint mean-and-precision predictions that a binary mean task cannot supply. The universal theory and its empirical validation remain unresolved.

\section{Three explicit mechanisms and truthful reporting}
A binary state $\theta\in\{0,1\}$ has prior log odds $L$. A symmetric binary signal $s\in\{-1,1\}$ has objective accuracy $q\in(1/2,1)$, so its log-likelihood ratio is $\ell=s\log(q/(1-q))$. Accuracy $q$, information price $c\geq0$ and desirability $v\in\{-1,1\}$ are randomized independently before signal realization. Conditional on each objective signal cell, the experiment has positive probability of either sign. The desirability label changes the benefit of believing the state to be one and, by assumption, does not change the state distribution or the incentive payment for accurate reporting.

Reports are log odds of the individual's subjective probability. Their measurement is either exact or adds error $\epsilon$ with $\E[\epsilon\mid L,\ell,c,v]=0$. This requires a valid elicitation model: for example, risk-neutral utility over a strictly proper score and no additional utility from the report itself. If beliefs and reports are separately motivated, the following identification applies to reports only and cannot automatically recover beliefs.

The fixed-misperception family $\mathcal M_\gamma$ uses perceived accuracy
$q_\gamma=\{1+\exp[-\gamma\log(q/(1-q))]\}^{-1}$, where $\gamma\geq0$ is invariant across tasks. It is a coherent subjective binary signal model, and yields
\[
 L^{M}=L+\gamma\ell.                                    \tag{1}
\]
The motivated family $\mathcal B_\eta$ first forms the correct Bayesian posterior $p_B$ and then chooses a belief $p\in[0,1]$ to maximize
\[
 \eta v p-\KL(\operatorname{Ber}(p)\Vert\operatorname{Ber}(p_B)),
 \qquad \eta\in\R.                                     \tag{2}
\]
The cost of departing from the Bayesian posterior is thus explicitly normalized. Finally the attention family $\mathcal A$ chooses whether to inspect before seeing $s$, with net material value $G(q,L)-\kappa c$, where $G$ is the expected gain from an explicitly specified downstream decision problem and $\kappa\geq0$. If inspection occurs, updating is Bayesian; otherwise the prior remains. With tie-breaking for inspection, this gives an indicator
$D=\Ind\{G(q,L)\geq\kappa c\}$ and
\[
 L^A=L+D\ell.                                          \tag{3}
\]
For a risk-neutral state-guessing task with prior one half and reward one for a correct guess, $G(q,0)=q-1/2$: with no signal the success probability is $1/2$, and with inspection it is $q$. The attention decision therefore has an explicit payoff basis.

\begin{lemma}[Motivated updating is an additive log-odds shift]
The unique optimizer of (2) is interior and satisfies
\[
 L^B=L+\ell+\eta v.                                    \tag{4}
\]
\end{lemma}
\begin{proof}
For $p\in(0,1)$, the derivative of the relative entropy in (2) is
$\log(p/(1-p))-\log(p_B/(1-p_B))$, and its second derivative is $1/[p(1-p)]>0$. The objective is strictly concave, with derivative tending to positive infinity at zero and negative infinity at one. Thus its unique first-order zero is the unique global maximizer, and rearranging that condition gives (4).
\end{proof}

\section{Factorial separation and its exact intersections}
Let $m(L,\ell,c,v)$ be the conditional mean report. Fix one prior and an accuracy with $\ell\in\{-x,x\}$, $x>0$, independently crossed with both values of $v$. Define
\[
 S(c,v)=\frac{m(L,x,c,v)-m(L,-x,c,v)}{2x},\qquad
 V(c,\ell)=\frac{m(L,\ell,c,1)-m(L,\ell,c,-1)}2.           \tag{5}
\]

\begin{proposition}[Separating restrictions]
Under the maintained reporting assumptions, $\mathcal M_\gamma$ implies $S=\gamma,V=0$, and $\mathcal B_\eta$ implies $S=1,V=\eta$. Equality of their complete conditional mean functions on this factorial design is possible if and only if $\gamma=1,\eta=0$. If attention has a distribution of $\kappa$ independent of $(q,c,v,s)$ and fixed prior, it implies
\[
 S(c,v)=\Pr\{G(q,L)\geq\kappa c\},\qquad V(c,\ell)=0.     \tag{6}
\]
For $G(q,L)\geq0$, the first quantity is nonincreasing in $c$. If it is strictly between zero and one in a treatment cell, exact report distributions distinguish it from every deterministic fixed-$\gamma$ model at that cell.
\end{proposition}
\begin{proof}
Substituting (1) and (4) into (5) proves the first two assertions. Equality of means forces $\gamma=1$ by comparing signal signs, then $\eta=0$ by comparing desirability signs; those parameter values indeed coincide. In the attention model $D$ is selected before the signal and the cost type is independent of signal and desirability. Hence the conditional mean is $L+\E[D\mid q,c,L]\ell$, giving (6). Increasing $c$ weakly shrinks $\{\kappa:\kappa c\leq G\}$, proving monotonicity. With exact reports and an interior inspection probability, conditional on a nonzero $\ell$ the attention law has two distinct atoms, $L$ and $L+\ell$, both with positive mass. A fixed-$\gamma$ model has one atom, $L+\gamma\ell$. A two-atom probability measure cannot equal a one-atom measure.
\end{proof}

The last claim can also hold with a \emph{known} independent additive error whose characteristic function never vanishes: equality of noisy laws would imply equality of latent laws after division of their characteristic functions, and uniqueness of characteristic functions then gives a contradiction. Unrestricted model-specific reporting errors would erase this conclusion. At the Bayesian boundary, attention with inspection probability one, misperception with $\gamma=1$, and motivation with $\eta=0$ are deliberately considered the same observational equivalence class. Attention with zero inspection also coincides with $\gamma=0$ if only posteriors are observed.

\section{A more fundamental nonidentification result}
The preceding separation depends on restrictions, especially fixed distortion and exclusion of desirability from information processing. Consider a broader misperception mechanism that erases the signal with probability $1-a(q,c,L)$, independently of its realization, and otherwise perceives it correctly. On erasure it perceives an uninformative null signal. The person knows the erasure process; hence null observation leaves the prior unchanged.

\begin{theorem}[Attention and erasure are observationally equivalent]
For every attention model (3) whose inspection probability conditional on $(q,c,L)$ is $a(q,c,L)$, there is a signal-erasure model with exactly the same law of reported beliefs under every assignment of $(q,c,v)$, including every finite sequence of such assignments with the same conditional history-dependent probabilities. Belief reports alone cannot identify which mechanism generated them.
\end{theorem}
\begin{proof}
Couple the two models using the same state, objective signal, reporting error, and Bernoulli variable $U$ having conditional success probability $a(q,c,L)$. Let attention inspect if $U=1$. Let the erasure model perceive the correct signal if $U=1$ and the null signal otherwise. In both models the posterior log odds are exactly $L+U\ell$ on every coupled realization. Thus their reports agree almost surely, proving equality of laws. For sequences, at each history draw the same next state/signal and the same Bernoulli variable conditional on that common history. Induction gives equality of the complete report paths. Randomization of the observed treatment coordinates does not alter the coupling argument.
\end{proof}

One additional intervention can separate these \emph{specified} models: externally supply a verified processed signal, setting $D=1$ in the attention model, while leaving the erasure probability unchanged in the competing perception model. If $0<a<1$ and $\ell\ne0$, the first model produces $L+\ell$ always and the second retains the two-atom mixture. This is a mathematical exclusion, not a claim that showing a signal on screen guarantees processing. A comprehension intervention must be validated to affect the two mechanisms differently; otherwise the equivalence persists. Likewise, a click record measures acquisition rather than comprehension unless a measurement model connects it to $D$.

\section{Cross-task mean and precision restrictions}
A second family uses a continuous state with prior $\theta\sim N(\mu_0,\tau_0^{-1})$ and observed signal $x=\theta+\varepsilon$, where $\varepsilon\sim N(0,\tau_s^{-1})$. Both positive precisions and their communication are known. Distributional reports are now elicited, with truthful reporting maintained. Extend fixed misperception by subjective signal precision $\gamma\tau_s$, $\gamma\geq0$, and extend motivated preferences by the objective
\[
 \eta v\E_Q[\theta]-\KL(Q\Vert P_B)                     \tag{7}
\]
over distributions $Q$ with the requisite finite expectations and relative entropy, where $P_B$ is the correct Gaussian posterior. Parameter sharing across binary and Gaussian tasks is a substantive restriction to test.

\begin{proposition}[Joint distribution predictions]
The Gaussian misperception and motivated-belief families imply respectively
\[
\begin{array}{c|c|c}
 &\text{posterior precision}&\text{posterior mean}\\ \hline
 \mathcal M_\gamma &\tau_0+\gamma\tau_s&
       \displaystyle\frac{\tau_0\mu_0+\gamma\tau_sx}{\tau_0+\gamma\tau_s}\\[6pt]
 \mathcal B_\eta&\tau_0+\tau_s&
       \displaystyle\frac{\tau_0\mu_0+\tau_sx+\eta v}{\tau_0+\tau_s}.
\end{array}                                                   \tag{8}
\]
In $\mathcal M_\gamma$, the mean's signal slope equals $1-\tau_0/T$, where $T$ is the reported posterior precision. In $\mathcal B_\eta$, desirability shifts the mean by $\eta v/T$ without changing precision.
\end{proposition}
\begin{proof}
Multiplying the Gaussian prior density by the Gaussian likelihood of precision $\gamma\tau_s$ and completing the square gives the first row, including $\gamma=0$ by ignoring the signal. For (7), define $Z=\E_{P_B}e^{\eta v\theta}<\infty$ and let $P^*$ have density $e^{\eta v\theta}/Z$ relative to $P_B$. Then, for every admissible $Q$,
\[
 \eta v\E_Q\theta-\KL(Q\Vert P_B)=\log Z-\KL(Q\Vert P^*).
\]
Nonnegativity of relative entropy, with equality only for equal distributions, proves that $P^*$ is the unique maximizer. Exponential tilting of the Gaussian changes its mean by $\eta v/(\tau_0+\tau_s)$ and leaves its precision fixed, as direct completion of the square shows. This gives the second row. Differentiating the first mean with respect to $x$ gives $\gamma\tau_s/(\tau_0+\gamma\tau_s)=1-\tau_0/T$; the remaining claim is visible in the second row.
\end{proof}

Thus a single parameter estimated from binary likelihood sensitivity predicts both Gaussian mean sensitivity and overprecision. Fixed likelihood distortion by itself does not create a special response to signals contradicting the prior: reversing the prior's sign leaves its slope unchanged. Failure on that confirmation task rejects this restricted unified model rather than justifying an additional unconstrained coefficient. Similarly, the motivated model in (7) cannot explain excessive reported precision without another primitive.

For a prespecified out-of-sample comparison on finite binned report outcomes, let a held-out treatment distribution be fixed before estimation and require every candidate forecast to assign positive probability to every outcome. Expected negative log score under the true conditional law $P$ differs from that of its own forecast by $\KL(P\Vert Q)\geq0$: expand $\sum_yP(y)\log[P(y)/Q(y)]$. Averaging over held-out cells retains nonnegativity. This provides a common scoring criterion, not a guarantee that a correctly fitted model wins in a finite sample. Confidence assessment still needs a sampling design and control of training selection.

The full research agenda therefore remains open: whether a small invariant mechanism family can jointly fit confirmation, precision, desirability, and prior-sensitivity behavior is an empirical question. The exact equivalence theorem identifies a limit of factorial report data, while (5) and (8) provide concrete restrictions that can fail and therefore make the proposed comparison informative.

\section{Completion Estimate}
60\%

This is a post hoc, heuristic estimate for the full catalogue target.
The attempt constructs competing belief mechanisms, proves factorial separation and attention-erasure equivalence, and derives shared binary and Gaussian mean-precision predictions. Full cross-task unification still requires validated reporting and processed-signal exclusions, confirmation and prior-sensitivity fit, and held-out scoring with sampling uncertainty.

\begin{thebibliography}{9}
\bibitem{rs} M. Rabin and J. L. Schrag (1999), \emph{First Impressions Matter: A Model of Confirmatory Bias}, Quarterly Journal of Economics 114(1), 37--82. Journal PDF mirrored by CREST, introduction and Section III inspected 9 October 2026. \url{https://crest.science/wp-content/uploads/2020/03/Rabin-and-Schrag-First-Impressions.pdf}. \url{https://doi.org/10.1162/003355399555945}.
\bibitem{bt} R. B\'enabou and J. Tirole (2002), \emph{Self-Confidence and Personal Motivation}, Quarterly Journal of Economics 117(3), 871--915. June 2001 author-hosted manuscript, introductory motivation inspected 9 October 2026; our KL model is not attributed to this paper. \url{https://www.princeton.edu/~rbenabou/papers/CONFQJE2.pdf}. \url{https://doi.org/10.1162/003355302760193913}.
\end{thebibliography}

\end{document}
